\section{The Memory Model}
\label{sec:model}

\stitle{Memories are values.}
A memory has a kind, a key, a body (any JSON value), the agent that wrote it, a
scope that says who may read it, an importance, and \code{derived\_from}. Its id
is a hash of all of that. Writing the same memory twice names the same memory,
and work that is redone and comes out the same restores what was there. A body
cannot change without its id changing, so tampering shows on load.

\stitle{State is a fold over a log.}
A record puts, restores, retracts, or marks stale one memory, or moves it between
tiers, inside a transaction. Status (\code{live}, \code{stale}, \code{retracted})
says whether a memory holds; tier (\code{short}, \code{mid}, \code{long},
\code{archive}) says where it sits. Neither is stored on the memory. Replaying the
log in order gives both, and \code{as\_of(lsn)} is a replay that stops early. A
transaction either commits all its records or none, and passes through guards
before the commit and observers after it. Rollback to an LSN writes the records
that return every memory to its status and tier there; the log still shows what
happened and what undid it.

\stitle{Provenance drives invalidation.}
\code{derived\_from} holds memory ids and outside references such as a document
version. Invalidating a reference marks stale every live memory reachable
through the reference index, in one transaction. A reference can also name a
\emph{read}: a memory built on ``no notes about this topic'' derives from the
query that returned nothing, and when a matching memory later goes live the work
built on the empty read goes stale. A retraction stales everything built on the
retracted memory, directly or through a read that included it.

\stitle{A declared structure.}
The core names no domain. A project declares its kinds as a schema: a JSON-schema
subset for each body, the scopes a kind may be written in, the kinds it must be
derived from, and a default tier. A commit guard rejects a write that is out of
line with its kind, for every writer in every process, and the schema is itself a
memory in the log, so \code{as\_of(lsn)} shows which structure held then. A kind
that \emph{requires} a premise (an answer derived from a question) declares
provenance together with structure.

\stitle{The cache and the database underneath.}
Short-term memory is a bounded cache in front of the log. A use counts as a hit,
a warm hit, or a miss; a missed live memory is paged back in; eviction policies
(FIFO, LRU, LFU, 2Q, ARC, a forgetting curve, a scored and a learned policy) are
pluggable. Compaction works at two levels: memory compaction folds finished
traces into long-term summaries and moves steps to the archive, and log
compaction folds each memory's history to the records that replay to the same
state, as an LSM tree rewrites sorted runs. Snapshots, hash-chain seals, a
retention policy, and a server that owns the log for agents in many processes
complete it.
